\section*{Findings}
The unconditional baseline comparisons show no statistically significant differences in female share, new hiring, internal transfers, office size, or average rank between offices drafted in 1944 and offices not drafted in 1944. Promotions are the exception in the unconditional comparison.

With section, position, and year fixed effects, future male draft intensity does not significantly predict female share, new hiring, promotion, office size, or average rank. It does positively predict the baseline count of internal transfers: one additional male drafted in 1944 is associated with 1.121 additional pre-war internal transfers ($p=0.025$). Thus the baseline tests do not establish complete balance on transfer intensity, even though unconditional transfer means are similar.

Future donor offices also show higher pre-war promotion counts ($p<0.001$). This estimate should be interpreted cautiously because the comparison contains 15 future donor offices and only one non-donor office.

Taken together with the era-by-stock analysis, the evidence suggests that future drafted offices differed somewhat in the \emph{amount} of pre-war internal mobility, but there is no detectable evidence that large- and small-1944-stock offices systematically drew transfers from different organizational source categories.
